\centerline{\bf \Huge Индукция в графах}

\begin{flushright}
        \scriptsize{Эпизод 18. Выбирая жизнь}
\end{flushright}


\begin{tcolorbox}[
    enhanced,
    breakable,
    colback=gray!10,
    colframe=black!70,
    boxrule=0.6pt,
    arc=2mm
]
В графах нельзя применять индукцию, делая переход от $k$ к $k+1$.
Схема индукции должна быть следующей: рассмотрим граф на $k+1$ вершине, удовлетворяющий условиям.
Удалим из него одну вершину так, чтобы граф, который получился, тоже удовлетворял условиям.
А в нем предположение индукции уже работает.

Самая сложная часть — понять, какую вершину можно удалить.
Для этого иногда нужно применять принцип крайнего — попробовать удалить вершину самой большой или самой маленькой степени и т.д.
\end{tcolorbox}

\problem{1}{Дан полный ориентированный граф на $n>2$ вершинах.
Докажите, что можно поменять ориентацию не более чем на одном ребре так, чтобы граф стал сильно связным.}

\problem{2}{В стране $n$ городов.
Между каждыми двумя из них проложена либо автомобильная, либо железная дорога.
Турист хочет объехать страну, побывав в каждом городе ровно один раз, и вернуться в город, с которого он начинал путешествие.
Докажите, что турист может выбрать город, с которого он начнет путешествие, и маршрут так, что ему придётся поменять вид транспорта не более одного раза.}

\problem{3}{На танцы девушек пришло больше, чем юношей.
Известно, что все девушки танцевали с разным числом юношей, а все юноши — с одинаковым числом девушек.
Сколько всего человек могло прийти на танцы?}

\problem{4}{В одной компании среди любых четырех человек есть знакомый с тремя остальными.
Докажите, что есть человек, который знает всех.}

\problem{5}{В графе $2n$ вершин и не меньше $n^2+1$ рёбер.
\textbf{(а)} Докажите, что в нём есть цикл длины три.
\textbf{(б)} Докажите, что в нем не менее $n$ циклов длины три.}

\problem{6}{Каждые два из $n$ городов соединены одной ориентированной синей или красной дорогой.
Докажите, что найдётся такой город, из которого в любой другой город можно проехать по одноцветному пути.}